\documentclass[12pt,a4paper,bibliography=totocnumbered,listof=totocnumbered]{scrartcl}
\usepackage[ngerman]{babel}
\usepackage[utf8]{inputenc}
\usepackage{amsmath}
\usepackage{amsfonts}
\usepackage{amssymb}
\usepackage{graphicx}
\usepackage{fancyhdr}
\usepackage{tabularx}
\usepackage{geometry}
\usepackage{setspace}
\usepackage[right]{eurosym}
\usepackage[printonlyused]{acronym}
\usepackage{subfig}
\usepackage{floatflt}
\usepackage[usenames,dvipsnames]{color}
\usepackage{colortbl}
\usepackage{paralist}
\usepackage{array}
\usepackage{titlesec}
\usepackage{parskip}
\usepackage[right]{eurosym}
\usepackage[subfigure,titles]{tocloft}
\usepackage[pdfpagelabels=true]{hyperref}
\usepackage[rgb]{xcolor}
\usepackage[obeyFinal, colorinlistoftodos, backgroundcolor=yellow, 
linecolor=yellow, textsize=tiny, textwidth=43]{todonotes}
\usepackage{letltxmacro}
\usepackage[autostyle, german=quotes]{csquotes}
\MakeOuterQuote{"}
\usepackage{biblatex}
\usepackage{enumitem}
\usepackage{etoolbox}
\usepackage{xparse}
\usepackage{import}
\usepackage[pagewise, right]{lineno}
\usepackage{ifdraft}
\usepackage{multirow}
\usepackage{placeins}

\usepackage{listings}
\lstset{basicstyle=\footnotesize, captionpos=b, breaklines=true, showstringspaces=false, tabsize=2, frame=lines, numbers=left, numberstyle=\tiny, xleftmargin=2em, framexleftmargin=2em}
\makeatletter
\def\l@lstlisting#1#2{\@dottedtocline{1}{0em}{1em}{\hspace{1,5em} Lst. #1}{#2}}
\makeatother

\geometry{a4paper, top=27mm, left=30mm, right=20mm, bottom=35mm, headsep=10mm, footskip=12mm}

\hypersetup{unicode=false, pdftoolbar=true, pdfmenubar=true, pdffitwindow=false, pdfstartview={FitH},
	pdftitle={Bachelorarbeit},
	pdfauthor={Rene Zarwel},
	pdfsubject={Microservices und technologische Heterogenität},
	pdfcreator={\LaTeX\ with package \flqq hyperref\frqq},
	pdfproducer={pdfTeX \the\pdftexversion.\pdftexrevision},
	pdfkeywords={Bachelorarbeit},
	pdfnewwindow=true,
	colorlinks=true,linkcolor=black,citecolor=black,filecolor=magenta,urlcolor=black}
\pdfinfo{/CreationDate (D:20110620133321)}

\setlist[description]{labelindent=1cm}
\setlist[itemize]{labelindent=1cm}
\setlist[enumerate]{labelindent=1cm}

\newcommand{\newparagraph}[1]{\paragraph{#1}~\\~\\}
\newcommand{\newsubparagraph}[1]{\subparagraph{#1}~\\}

\newcommand{\acfootnote}[1]{\acs{#1}\footnote{\acf{#1}}}

% Nice and big quote
% #1 Quote
% #2 Name of person to quote
\newcommand{\bigquote}[2]{
	\vspace*{1\baselineskip}
	\begin{center}
		\begin{minipage}{0.8\textwidth}
			\enquote{\textit{#1}}\par
			- #2
		\end{minipage}
	\end{center}
	\vspace*{1\baselineskip}
}
%-----------------

%Todo Commands
\newcounter{todocounter}

\LetLtxMacro{\todoold}{\todo}
\renewcommand{\todo}[2][]{\stepcounter{todocounter}\todoold[#1]{TODO \thetodocounter: #2}}
\newcommand{\todoInline}[1]{\todo[inline]{#1}}
\newcommand{\todoFigure}[1]{\stepcounter{todocounter}\missingfigure[figcolor=white]{TODO \thetodocounter: #1}}
% ----------------

%Cite Redefinition - Auto Space on citation + Prevent forever alone citation
\renewcommand{\multicitedelim}{\addsemicolon\space}
\LetLtxMacro{\citeold}{\cite}
\renewcommand{\cite}[2][]{~\citeold[#1]{#2}}
\LetLtxMacro{\citesold}{\cites}
\DeclareDocumentCommand{\cites}{o m o m}{
	\IfNoValueTF{#1}{
		\IfNoValueTF{#3}{
			~\citesold{#2}{#4}
		}{
			~\citesold{#2}[#3]{#4}
		}
	}{
		\IfNoValueTF{#3}{
			~\citesold[#1]{#2}{#4}
		}{
			~\citesold[#1]{#2}[#3]{#4}
		}
	}
}
%----------------------

%Image Command - Optional Image otherwise TODO
% #1 IMAGE File - optional (TODO if not present)
% #2 Image options - optional
% #3 LabelName
% #4 TOC Text
% #5 Captiontext
\DeclareDocumentCommand{\image}{ o O{width=\linewidth} m m m}{
	\begin{minipage}{\linewidth}
		\vspace*{0.5cm}
		\centering
		\IfNoValueTF{#1}{
			\todoFigure{#4}
		}{
			\includegraphics[#2]{Bilder/#1}
		}	
		\captionof{figure}[#4]{#5}
		\label{#3}
		\vspace*{0.5cm}
\end{minipage}
}
%-------------

\bibliography{bibo}

\begin{document}

\titlespacing{\section}{0pt}{12pt plus 4pt minus 2pt}{-6pt plus 2pt minus 2pt}

% Kopf- und Fusszeile
\renewcommand{\sectionmark}[1]{\markright{#1}}
\renewcommand{\leftmark}{\rightmark}
\pagestyle{fancy}
\lhead{}
\chead{}
\rhead{\thesection\space\contentsname}
\lfoot{Microservices und technologische Heterogenität}
\cfoot{}
\rfoot{\ \linebreak Seite \thepage}
\renewcommand{\headrulewidth}{0.4pt}
\renewcommand{\footrulewidth}{0.4pt}

% Vorspann
\renewcommand{\thesection}{\Roman{section}}
\renewcommand{\theHsection}{\Roman{section}}
\pagenumbering{Roman}

% ----------------------------------------------------------------------------------------------------------
% Titelseite
% ----------------------------------------------------------------------------------------------------------
\thispagestyle{empty}
\begin{center}
	\includegraphics[scale=0.5]{Bilder/Hochschule_Muenchen_Logo.png}\\
	\vspace*{2cm}
	\Large
	\textbf{Fakultät für Informatik und Mathematik 07}\\
	\vspace*{2cm}
	\Huge
	\textbf{Bachelorarbeit}\\
	\vspace*{0.5cm}
	\large
	über das Thema\\
	\vspace*{1cm}
	\textbf{Microservices und technologische Heterogenität}\\
	Entwicklung einer sprachunabhängigen Microservice Framework Evaluationsmethode\\
	\vspace*{2cm}
	
	\vfill
	\normalsize
	\newcolumntype{x}[1]{>{\raggedleft\arraybackslash\hspace{0pt}}p{#1}}
	\begin{tabular}{x{6cm}p{7.5cm}}
		\rule{0mm}{5ex}\textbf{Autor:} & René Zarwel\newline zarwel@hm.edu \\ 
		\rule{0mm}{5ex}\textbf{Prüfer:} & Prof. Dr. Hammerschall \\ 
		\rule{0mm}{5ex}\textbf{Abgabedatum:} & XX.XX.XX\\ 
	\end{tabular} 
\end{center}
\pagebreak

% ----------------------------------------------------------------------------------------------------------
% Abstract
% ----------------------------------------------------------------------------------------------------------
\setcounter{page}{1}
\onehalfspacing
\titlespacing{\section}{0pt}{12pt plus 4pt minus 2pt}{2pt plus 2pt minus 2pt}
\rhead{KURZFASSUNG}
\section{Kurzfassung}
Insert Abstract Here

\vspace{-1,2em}
\titlespacing{\section}{0pt}{12pt plus 4pt minus 2pt}{-6pt plus 2pt minus 2pt}
\section*{Abstract}
Insert Abstract Here
\pagebreak

% ----------------------------------------------------------------------------------------------------------
% Verzeichnisse
% ----------------------------------------------------------------------------------------------------------
% TODO Typ vor Nummer
\renewcommand{\cfttabpresnum}{Tab. }
\renewcommand{\cftfigpresnum}{Abb. }
\settowidth{\cfttabnumwidth}{Abb. 10\quad}
\settowidth{\cftfignumwidth}{Abb. 10\quad}

\titlespacing{\section}{0pt}{12pt plus 4pt minus 2pt}{2pt plus 2pt minus 2pt}
\singlespacing
\rhead{INHALTSVERZEICHNIS}
\renewcommand{\contentsname}{II Inhaltsverzeichnis}
\phantomsection
\addcontentsline{toc}{section}{\texorpdfstring{II \hspace{0.35em}Inhaltsverzeichnis}{Inhaltsverzeichnis}}
\addtocounter{section}{1}
\tableofcontents
\pagebreak
\rhead{VERZEICHNISSE}
\listoffigures
\pagebreak
\listoftables
%\pagebreak
\renewcommand{\lstlistlistingname}{Listing-Verzeichnis}
{\labelsep2cm\lstlistoflistings}
\pagebreak
% ----------------------------------------------------------------------------------------------------------
% Abkürzungen
% ----------------------------------------------------------------------------------------------------------
\section{Abkürzungsverzeichnis}
\begin{acronym}[ATAM] % längste Abkürzung steht in eckigen Klammern
	\setlength{\itemsep}{-\parsep} % geringerer Zeilenabstand
	\acro{SEI}{Software Engineering Institute}
	\acro{SAAM}{Software Architecture Analysis Method}
	\acro{ATAM}{Architecture Tradeoff Analysis Method}
	\acro{SAEM}{Software Architecture Evaluation Model}
	\acro{GQM}{Goal Question Metrik}
	\acro{MFEM}{Microservice Framework Evaluation Method}
	\acro{HATEOAS}{Hypermedia As The Engine Of Application State}
	\acro{REST}{Representational State Transfer}
	\acro{XML}{Extensible Markup Language}
	\acro{JSON}{JavaScript Object Notation}
	\acro{CW}{Cognitive Walkthrough}
	\acro{JWT}{JSON Web Token}
	\acro{API}{Programmierschnittstelle}
\end{acronym}
\newpage

% ----------------------------------------------------------------------------------------------------------
% Inhalt
% ----------------------------------------------------------------------------------------------------------
% Abstände Überschrift
\titlespacing{\section}{0pt}{12pt plus 4pt minus 2pt}{-6pt plus 2pt minus 2pt}
\titlespacing{\subsection}{0pt}{12pt plus 4pt minus 2pt}{-6pt plus 2pt minus 2pt}
\titlespacing{\subsubsection}{0pt}{12pt plus 4pt minus 2pt}{-6pt plus 2pt minus 2pt}

% Kopfzeile
\renewcommand{\sectionmark}[1]{\markright{#1}}
\renewcommand{\subsectionmark}[1]{}
\renewcommand{\subsubsectionmark}[1]{}
\lhead{Kapitel \thesection}
\rhead{\rightmark}

\onehalfspacing
\renewcommand{\thesection}{\arabic{section}}
\renewcommand{\theHsection}{\arabic{section}}
\setcounter{section}{0}
\pagenumbering{arabic}
\setcounter{page}{1}

%
\ifoptionfinal{}{\linenumbers}
 ----------------------------------------------------------------------------------------------------------
% Einleitung und Motivation
% ----------------------------------------------------------------------------------------------------------
\subimport{Einleitung/}{Einleitung}
\pagebreak
% ----------------------------------------------------------------------------------------------------------
% Microservices
% ----------------------------------------------------------------------------------------------------------
\subimport{Grundlagen/}{Grundlagen}
\pagebreak

% ----------------------------------------------------------------------------------------------------------
% Konzeption Methode
% ----------------------------------------------------------------------------------------------------------
\subimport{Konzeption/}{Konzeption}

\pagebreak

% ----------------------------------------------------------------------------------------------------------
% Evaluation an Beispielen
% ----------------------------------------------------------------------------------------------------------
\subimport{Evaluation/}{Evaluation}

\pagebreak

% ----------------------------------------------------------------------------------------------------------
% Zusammenfassung und Ausblick
% ----------------------------------------------------------------------------------------------------------
\subimport{Zusammenfassung/}{Zusammenfassung}

\pagebreak
\nolinenumbers
% ----------------------------------------------------------------------------------------------------------
% Literatur
% ----------------------------------------------------------------------------------------------------------
\renewcommand\refname{Quellenverzeichnis}

\printbibliography
\pagebreak

% ----------------------------------------------------------------------------------------------------------
% Todos
% ----------------------------------------------------------------------------------------------------------
\todototoc
\listoftodos

% ----------------------------------------------------------------------------------------------------------
% Anhang
% ----------------------------------------------------------------------------------------------------------
\pagenumbering{Roman}
\setcounter{page}{1}
\lhead{Anhang \thesection}

\begin{appendix}
\section*{Anhang}
\phantomsection
\addcontentsline{toc}{section}{Anhang}
\addtocontents{toc}{\vspace{-0.5em}}
\section{\ac{MFEM} - Vollständiger Quality Utility Tree}
\todoInline{Quality Utility Tree vollständig mit allen Kategorien (mit GQM)}

\section{CODE}
\todoInline{Nothing Here}

\end{appendix}
\end{document}
